\subsection{DIVIDED POWERS IN A BIGEBRA}

Let A be a unital k-algebra, \(x \in A\) , \(d \in k\) , \(n \in N\) . Put

\[
x ^ {(n, d)} = \frac {x (x - d) \dots (x - d (n - 1))}{n !} = \prod_ {i = 0} ^ {n - 1} (x - i d) / (i + 1).\tag{1}
\]

In particular, \(x^{(0,d)} = 1\) , \(x^{(1,d)} = x\) . We agree that \(x^{(n,d)} = 0\) for \(n\) an integer \(< 0\) . Put

\[
x ^ {(n)} = x ^ {(n, 0)} = \frac {x ^ {n}}{n !}\tag{2}
\]

\[
\binom {x} {n} = x ^ {(n, 1)} = \frac {x (x - 1) \ldots (x - n + 1)}{n !}.\tag{3}
\]

PROPOSITION 1. Let A be a bigebra, with coproduct \(c\) , and \(x\) a primitive element (Chap. II, §1, no. 2) of A. Then

\[
c (x ^ {(n, d)}) = \sum_ {p \in \mathbf {N}, q \in \mathbf {N}, p + q = n} x ^ {(p, d)} \otimes x ^ {(q, d)}.\tag{4}
\]

The proposition is trivial for \(n \leq 0\) . We argue by induction on \(n\) . If formula (4) is true for \(n\) , then

\[
\begin{array}{l} (n + 1) c (x ^ {(n + 1, d)}) = c (x - d n) c (x ^ {(n, d)}) \\ \quad = (x \otimes 1 + 1 \otimes x - d n   1 \otimes 1) c (x ^ {(n, d)}) \\ \quad = \sum_ {p + q = n} [ x x ^ {(p, d)} \otimes x ^ {(q, d)} + x ^ {(p, d)} \otimes x x ^ {(q, d)} - (p + q) d x ^ {(p, d)} \otimes x ^ {(q, d)} ] \\ \quad = \sum_ {p + q = n} (x - p d) x ^ {(p, d)} \otimes x ^ {(q, d)} + \sum_ {p + q = n} x ^ {(p, d)} \otimes (x - q d) x ^ {(q, d)} \\ \quad = \sum_ {p + q = n} (p + 1) x ^ {(p + 1, d)} \otimes x ^ {(q, d)} + \sum_ {p + q = n} (q + 1) x ^ {(p, d)} \otimes x ^ {(q + 1, d)} \\ \quad = \sum_ {r + s = n + 1} r x ^ {(r, d)} \otimes x ^ {(s, d)} + \sum_ {r + s = n + 1} s x ^ {(r, d)} \otimes x ^ {(s, d)} \\ \quad = (n + 1) \sum_ {r + s = n + 1} x ^ {(r, d)} \otimes x ^ {(s, d)}, \end{array}
\]

hence formula (4) for \(n + 1\) .

